\documentclass[conference]{IEEEtran}
\usepackage{graphicx,booktabs,amsmath,amssymb,cite,url}
\graphicspath{{../figures/}}
% AUTO-GENERATED by scripts/make_report_assets.py. Do not edit by hand.
\newcommand{\congOffpeak}{5.8}
\newcommand{\congPeak}{26.9}
\newcommand{\corrEightHop}{0.14}
\newcommand{\corrOneHop}{0.45}
\newcommand{\corrTwoHop}{0.25}
\newcommand{\gapAllFifteen}{0.059}
\newcommand{\gapAllFifteenHi}{0.069}
\newcommand{\gapAllFifteenLo}{0.051}
\newcommand{\gapAllSixty}{0.051}
\newcommand{\gapAllSixtyHi}{0.072}
\newcommand{\gapAllSixtyLo}{0.033}
\newcommand{\gapAllThirty}{0.073}
\newcommand{\gapAllThirtyHi}{0.090}
\newcommand{\gapAllThirtyLo}{0.059}
\newcommand{\gapCongestedFifteen}{0.336}
\newcommand{\gapCongestedFifteenHi}{0.392}
\newcommand{\gapCongestedFifteenLo}{0.282}
\newcommand{\gapCongestedSixty}{0.458}
\newcommand{\gapCongestedSixtyHi}{0.570}
\newcommand{\gapCongestedSixtyLo}{0.343}
\newcommand{\gapCongestedThirty}{0.559}
\newcommand{\gapCongestedThirtyHi}{0.670}
\newcommand{\gapCongestedThirtyLo}{0.452}
\newcommand{\gapFreeflowFifteen}{0.020}
\newcommand{\gapFreeflowFifteenHi}{0.024}
\newcommand{\gapFreeflowFifteenLo}{0.017}
\newcommand{\gapFreeflowSixty}{-0.007}
\newcommand{\gapFreeflowSixtyHi}{0.004}
\newcommand{\gapFreeflowSixtyLo}{-0.018}
\newcommand{\gapFreeflowThirty}{0.005}
\newcommand{\gapFreeflowThirtyHi}{0.012}
\newcommand{\gapFreeflowThirtyLo}{-0.003}
\newcommand{\gapOffpeakFifteen}{0.038}
\newcommand{\gapOffpeakFifteenHi}{0.048}
\newcommand{\gapOffpeakFifteenLo}{0.031}
\newcommand{\gapOffpeakSixty}{0.032}
\newcommand{\gapOffpeakSixtyHi}{0.055}
\newcommand{\gapOffpeakSixtyLo}{0.012}
\newcommand{\gapOffpeakThirty}{0.048}
\newcommand{\gapOffpeakThirtyHi}{0.065}
\newcommand{\gapOffpeakThirtyLo}{0.034}
\newcommand{\gapPeakFifteen}{0.117}
\newcommand{\gapPeakFifteenHi}{0.135}
\newcommand{\gapPeakFifteenLo}{0.101}
\newcommand{\gapPeakSixty}{0.112}
\newcommand{\gapPeakSixtyHi}{0.143}
\newcommand{\gapPeakSixtyLo}{0.081}
\newcommand{\gapPeakThirty}{0.152}
\newcommand{\gapPeakThirtyHi}{0.182}
\newcommand{\gapPeakThirtyLo}{0.126}
\newcommand{\gapWeekendFifteen}{0.036}
\newcommand{\gapWeekendFifteenHi}{0.057}
\newcommand{\gapWeekendFifteenLo}{0.025}
\newcommand{\gapWeekendSixty}{0.020}
\newcommand{\gapWeekendSixtyHi}{0.041}
\newcommand{\gapWeekendSixtyLo}{0.000}
\newcommand{\gapWeekendThirty}{0.037}
\newcommand{\gapWeekendThirtyHi}{0.057}
\newcommand{\gapWeekendThirtyLo}{0.024}
\newcommand{\graphDiameter}{17}
\newcommand{\graphEdges}{1515}
\newcommand{\graphNodes}{207}
\newcommand{\graphSCC}{195}
\newcommand{\maeHistoricalAverageFifteen}{4.19}
\newcommand{\maeHistoricalAverageFifteenCongested}{12.90}
\newcommand{\maeHistoricalAverageFifteenFreeflow}{2.96}
\newcommand{\maeHistoricalAverageFifteenOffpeak}{3.52}
\newcommand{\maeHistoricalAverageFifteenPeak}{6.04}
\newcommand{\maeHistoricalAverageFifteenPeakAm}{4.90}
\newcommand{\maeHistoricalAverageFifteenPeakPm}{7.13}
\newcommand{\maeHistoricalAverageFifteenWeekend}{3.50}
\newcommand{\maeHistoricalAverageSixty}{4.19}
\newcommand{\maeHistoricalAverageSixtyCongested}{12.90}
\newcommand{\maeHistoricalAverageSixtyFreeflow}{2.96}
\newcommand{\maeHistoricalAverageSixtyOffpeak}{3.51}
\newcommand{\maeHistoricalAverageSixtyPeak}{6.06}
\newcommand{\maeHistoricalAverageSixtyPeakAm}{4.93}
\newcommand{\maeHistoricalAverageSixtyPeakPm}{7.13}
\newcommand{\maeHistoricalAverageSixtyWeekend}{3.50}
\newcommand{\maeHistoricalAverageThirty}{4.19}
\newcommand{\maeHistoricalAverageThirtyCongested}{12.90}
\newcommand{\maeHistoricalAverageThirtyFreeflow}{2.96}
\newcommand{\maeHistoricalAverageThirtyOffpeak}{3.52}
\newcommand{\maeHistoricalAverageThirtyPeak}{6.05}
\newcommand{\maeHistoricalAverageThirtyPeakAm}{4.91}
\newcommand{\maeHistoricalAverageThirtyPeakPm}{7.13}
\newcommand{\maeHistoricalAverageThirtyWeekend}{3.50}
\newcommand{\maePersistenceFifteen}{3.49}
\newcommand{\maePersistenceFifteenCongested}{7.31}
\newcommand{\maePersistenceFifteenFreeflow}{2.96}
\newcommand{\maePersistenceFifteenOffpeak}{3.29}
\newcommand{\maePersistenceFifteenPeak}{4.39}
\newcommand{\maePersistenceFifteenPeakAm}{3.75}
\newcommand{\maePersistenceFifteenPeakPm}{4.99}
\newcommand{\maePersistenceFifteenWeekend}{2.86}
\newcommand{\maePersistenceSixty}{5.37}
\newcommand{\maePersistenceSixtyCongested}{13.33}
\newcommand{\maePersistenceSixtyFreeflow}{4.24}
\newcommand{\maePersistenceSixtyOffpeak}{4.85}
\newcommand{\maePersistenceSixtyPeak}{7.45}
\newcommand{\maePersistenceSixtyPeakAm}{6.93}
\newcommand{\maePersistenceSixtyPeakPm}{7.95}
\newcommand{\maePersistenceSixtyWeekend}{3.98}
\newcommand{\maePersistenceThirty}{4.20}
\newcommand{\maePersistenceThirtyCongested}{9.63}
\newcommand{\maePersistenceThirtyFreeflow}{3.43}
\newcommand{\maePersistenceThirtyOffpeak}{3.84}
\newcommand{\maePersistenceThirtyPeak}{5.58}
\newcommand{\maePersistenceThirtyPeakAm}{4.96}
\newcommand{\maePersistenceThirtyPeakPm}{6.17}
\newcommand{\maePersistenceThirtyWeekend}{3.31}
\newcommand{\maeRfMultioutputFifteen}{3.14}
\newcommand{\maeRfMultioutputFifteenCongested}{7.68}
\newcommand{\maeRfMultioutputFifteenFreeflow}{2.49}
\newcommand{\maeRfMultioutputFifteenOffpeak}{2.80}
\newcommand{\maeRfMultioutputFifteenPeak}{4.21}
\newcommand{\maeRfMultioutputFifteenPeakAm}{3.59}
\newcommand{\maeRfMultioutputFifteenPeakPm}{4.81}
\newcommand{\maeRfMultioutputFifteenWeekend}{2.59}
\newcommand{\maeRfMultioutputSixty}{4.73}
\newcommand{\maeRfMultioutputSixtyCongested}{14.05}
\newcommand{\maeRfMultioutputSixtyFreeflow}{3.42}
\newcommand{\maeRfMultioutputSixtyOffpeak}{3.82}
\newcommand{\maeRfMultioutputSixtyPeak}{7.24}
\newcommand{\maeRfMultioutputSixtyPeakAm}{6.52}
\newcommand{\maeRfMultioutputSixtyPeakPm}{7.91}
\newcommand{\maeRfMultioutputSixtyWeekend}{3.80}
\newcommand{\maeRfMultioutputThirty}{3.80}
\newcommand{\maeRfMultioutputThirtyCongested}{10.29}
\newcommand{\maeRfMultioutputThirtyFreeflow}{2.88}
\newcommand{\maeRfMultioutputThirtyOffpeak}{3.24}
\newcommand{\maeRfMultioutputThirtyPeak}{5.44}
\newcommand{\maeRfMultioutputThirtyPeakAm}{4.79}
\newcommand{\maeRfMultioutputThirtyPeakPm}{6.05}
\newcommand{\maeRfMultioutputThirtyWeekend}{3.08}
\newcommand{\maeSeasonalNaiveFifteen}{5.50}
\newcommand{\maeSeasonalNaiveFifteenCongested}{14.79}
\newcommand{\maeSeasonalNaiveFifteenFreeflow}{4.19}
\newcommand{\maeSeasonalNaiveFifteenOffpeak}{4.38}
\newcommand{\maeSeasonalNaiveFifteenPeak}{7.45}
\newcommand{\maeSeasonalNaiveFifteenPeakAm}{6.14}
\newcommand{\maeSeasonalNaiveFifteenPeakPm}{8.70}
\newcommand{\maeSeasonalNaiveFifteenWeekend}{5.83}
\newcommand{\maeSeasonalNaiveSixty}{5.50}
\newcommand{\maeSeasonalNaiveSixtyCongested}{14.80}
\newcommand{\maeSeasonalNaiveSixtyFreeflow}{4.19}
\newcommand{\maeSeasonalNaiveSixtyOffpeak}{4.37}
\newcommand{\maeSeasonalNaiveSixtyPeak}{7.48}
\newcommand{\maeSeasonalNaiveSixtyPeakAm}{6.18}
\newcommand{\maeSeasonalNaiveSixtyPeakPm}{8.70}
\newcommand{\maeSeasonalNaiveSixtyWeekend}{5.83}
\newcommand{\maeSeasonalNaiveThirty}{5.50}
\newcommand{\maeSeasonalNaiveThirtyCongested}{14.79}
\newcommand{\maeSeasonalNaiveThirtyFreeflow}{4.19}
\newcommand{\maeSeasonalNaiveThirtyOffpeak}{4.38}
\newcommand{\maeSeasonalNaiveThirtyPeak}{7.46}
\newcommand{\maeSeasonalNaiveThirtyPeakAm}{6.15}
\newcommand{\maeSeasonalNaiveThirtyPeakPm}{8.70}
\newcommand{\maeSeasonalNaiveThirtyWeekend}{5.83}
\newcommand{\maeXgbGlobalFifteen}{2.88}
\newcommand{\maeXgbGlobalFifteenCongested}{7.51}
\newcommand{\maeXgbGlobalFifteenFreeflow}{2.23}
\newcommand{\maeXgbGlobalFifteenOffpeak}{2.58}
\newcommand{\maeXgbGlobalFifteenPeak}{3.87}
\newcommand{\maeXgbGlobalFifteenPeakAm}{3.21}
\newcommand{\maeXgbGlobalFifteenPeakPm}{4.49}
\newcommand{\maeXgbGlobalFifteenWeekend}{2.36}
\newcommand{\maeXgbGlobalSixty}{3.78}
\newcommand{\maeXgbGlobalSixtyCongested}{13.03}
\newcommand{\maeXgbGlobalSixtyFreeflow}{2.48}
\newcommand{\maeXgbGlobalSixtyOffpeak}{3.11}
\newcommand{\maeXgbGlobalSixtyPeak}{5.65}
\newcommand{\maeXgbGlobalSixtyPeakAm}{4.67}
\newcommand{\maeXgbGlobalSixtyPeakPm}{6.58}
\newcommand{\maeXgbGlobalSixtyWeekend}{3.05}
\newcommand{\maeXgbGlobalThirty}{3.33}
\newcommand{\maeXgbGlobalThirtyCongested}{10.15}
\newcommand{\maeXgbGlobalThirtyFreeflow}{2.37}
\newcommand{\maeXgbGlobalThirtyOffpeak}{2.86}
\newcommand{\maeXgbGlobalThirtyPeak}{4.73}
\newcommand{\maeXgbGlobalThirtyPeakAm}{3.99}
\newcommand{\maeXgbGlobalThirtyPeakPm}{5.44}
\newcommand{\maeXgbGlobalThirtyWeekend}{2.70}
\newcommand{\maeXgbPerSegmentFifteen}{2.86}
\newcommand{\maeXgbPerSegmentFifteenCongested}{7.63}
\newcommand{\maeXgbPerSegmentFifteenFreeflow}{2.19}
\newcommand{\maeXgbPerSegmentFifteenOffpeak}{2.57}
\newcommand{\maeXgbPerSegmentFifteenPeak}{3.78}
\newcommand{\maeXgbPerSegmentFifteenPeakAm}{3.06}
\newcommand{\maeXgbPerSegmentFifteenPeakPm}{4.48}
\newcommand{\maeXgbPerSegmentFifteenWeekend}{2.38}
\newcommand{\maeXgbPerSegmentSixty}{3.68}
\newcommand{\maeXgbPerSegmentSixtyCongested}{12.15}
\newcommand{\maeXgbPerSegmentSixtyFreeflow}{2.49}
\newcommand{\maeXgbPerSegmentSixtyOffpeak}{3.11}
\newcommand{\maeXgbPerSegmentSixtyPeak}{5.26}
\newcommand{\maeXgbPerSegmentSixtyPeakAm}{4.19}
\newcommand{\maeXgbPerSegmentSixtyPeakPm}{6.28}
\newcommand{\maeXgbPerSegmentSixtyWeekend}{3.09}
\newcommand{\maeXgbPerSegmentThirty}{3.28}
\newcommand{\maeXgbPerSegmentThirtyCongested}{9.91}
\newcommand{\maeXgbPerSegmentThirtyFreeflow}{2.34}
\newcommand{\maeXgbPerSegmentThirtyOffpeak}{2.86}
\newcommand{\maeXgbPerSegmentThirtyPeak}{4.51}
\newcommand{\maeXgbPerSegmentThirtyPeakAm}{3.63}
\newcommand{\maeXgbPerSegmentThirtyPeakPm}{5.35}
\newcommand{\maeXgbPerSegmentThirtyWeekend}{2.73}
\newcommand{\maeXgbSpatialFifteen}{2.82}
\newcommand{\maeXgbSpatialFifteenCongested}{7.18}
\newcommand{\maeXgbSpatialFifteenFreeflow}{2.21}
\newcommand{\maeXgbSpatialFifteenOffpeak}{2.54}
\newcommand{\maeXgbSpatialFifteenPeak}{3.75}
\newcommand{\maeXgbSpatialFifteenPeakAm}{3.14}
\newcommand{\maeXgbSpatialFifteenPeakPm}{4.33}
\newcommand{\maeXgbSpatialFifteenWeekend}{2.32}
\newcommand{\maeXgbSpatialSixty}{3.73}
\newcommand{\maeXgbSpatialSixtyCongested}{12.58}
\newcommand{\maeXgbSpatialSixtyFreeflow}{2.48}
\newcommand{\maeXgbSpatialSixtyOffpeak}{3.08}
\newcommand{\maeXgbSpatialSixtyPeak}{5.54}
\newcommand{\maeXgbSpatialSixtyPeakAm}{4.65}
\newcommand{\maeXgbSpatialSixtyPeakPm}{6.38}
\newcommand{\maeXgbSpatialSixtyWeekend}{3.03}
\newcommand{\maeXgbSpatialThirty}{3.26}
\newcommand{\maeXgbSpatialThirtyCongested}{9.59}
\newcommand{\maeXgbSpatialThirtyFreeflow}{2.36}
\newcommand{\maeXgbSpatialThirtyOffpeak}{2.82}
\newcommand{\maeXgbSpatialThirtyPeak}{4.58}
\newcommand{\maeXgbSpatialThirtyPeakAm}{3.88}
\newcommand{\maeXgbSpatialThirtyPeakPm}{5.24}
\newcommand{\maeXgbSpatialThirtyWeekend}{2.66}
\newcommand{\mapeHistoricalAverageFifteen}{13.1}
\newcommand{\mapeHistoricalAverageSixty}{13.1}
\newcommand{\mapeHistoricalAverageThirty}{13.1}
\newcommand{\mapePersistenceFifteen}{8.6}
\newcommand{\mapePersistenceSixty}{14.5}
\newcommand{\mapePersistenceThirty}{10.8}
\newcommand{\mapeRfMultioutputFifteen}{8.5}
\newcommand{\mapeRfMultioutputSixty}{14.3}
\newcommand{\mapeRfMultioutputThirty}{10.9}
\newcommand{\mapeSeasonalNaiveFifteen}{15.6}
\newcommand{\mapeSeasonalNaiveSixty}{15.6}
\newcommand{\mapeSeasonalNaiveThirty}{15.6}
\newcommand{\mapeXgbGlobalFifteen}{7.9}
\newcommand{\mapeXgbGlobalSixty}{12.1}
\newcommand{\mapeXgbGlobalThirty}{9.9}
\newcommand{\mapeXgbPerSegmentFifteen}{7.9}
\newcommand{\mapeXgbPerSegmentSixty}{11.7}
\newcommand{\mapeXgbPerSegmentThirty}{9.8}
\newcommand{\mapeXgbSpatialFifteen}{7.6}
\newcommand{\mapeXgbSpatialSixty}{11.8}
\newcommand{\mapeXgbSpatialThirty}{9.5}
\newcommand{\meanSpeed}{58.5}
\newcommand{\missingRate}{8.1}
\newcommand{\missingTest}{12.1}
\newcommand{\missingTrain}{7.1}
\newcommand{\missingVal}{7.0}
\newcommand{\networkOutages}{2148}
\newcommand{\rmseHistoricalAverageFifteen}{7.86}
\newcommand{\rmseHistoricalAverageSixty}{7.86}
\newcommand{\rmseHistoricalAverageThirty}{7.86}
\newcommand{\rmsePersistenceFifteen}{6.39}
\newcommand{\rmsePersistenceSixty}{10.30}
\newcommand{\rmsePersistenceThirty}{8.02}
\newcommand{\rmseRfMultioutputFifteen}{5.70}
\newcommand{\rmseRfMultioutputSixty}{8.25}
\newcommand{\rmseRfMultioutputThirty}{6.89}
\newcommand{\rmseSeasonalNaiveFifteen}{10.71}
\newcommand{\rmseSeasonalNaiveSixty}{10.71}
\newcommand{\rmseSeasonalNaiveThirty}{10.71}
\newcommand{\rmseXgbGlobalFifteen}{5.69}
\newcommand{\rmseXgbGlobalSixty}{7.76}
\newcommand{\rmseXgbGlobalThirty}{6.80}
\newcommand{\rmseXgbPerSegmentFifteen}{5.60}
\newcommand{\rmseXgbPerSegmentSixty}{7.52}
\newcommand{\rmseXgbPerSegmentThirty}{6.62}
\newcommand{\rmseXgbSpatialFifteen}{5.46}
\newcommand{\rmseXgbSpatialSixty}{7.54}
\newcommand{\rmseXgbSpatialThirty}{6.54}
\newcommand{\splitTestEnd}{2012-06-27}
\newcommand{\splitTestStart}{2012-06-04}
\newcommand{\splitTrainEnd}{2012-05-23}
\newcommand{\splitTrainStart}{2012-03-01}
\newcommand{\splitValEnd}{2012-06-04}
\newcommand{\splitValStart}{2012-05-23}
\newcommand{\xcorrAllNeg}{30}
\newcommand{\xcorrAllPos}{38}
\newcommand{\xcorrAllZero}{32}
\newcommand{\xcorrPairs}{1228}
\newcommand{\xcorrPeakNeg}{38}
\newcommand{\xcorrPeakPos}{50}
\newcommand{\xcorrPeakZero}{12}

\begin{document}
\title{Traffic Speed Forecasting on METR-LA: Segment-Level Machine Learning and the Case for Graph Models (Phase 1)}
\author{\IEEEauthorblockN{Pranav Singh, Kunal, Himanshu, Kanchan, and Vaibhav}}
\maketitle

\begin{abstract}
We forecast freeway speed 5--60 minutes ahead on the METR-LA sensor network (\graphNodes{} loop detectors, 4 months) and ask whether information from neighbouring sensors helps beyond each sensor's own history.
Using a strict chronological protocol, masked metrics and train-only preprocessing, we compare naive baselines, per-segment and pooled XGBoost, a multi-output Random Forest and an XGBoost model with hand-engineered neighbour features.
The best models reduce MAE at 15 minutes from \maePersistenceFifteen{} mph (persistence) to \maeXgbSpatialFifteen{} mph, and at 60 minutes from \maePersistenceSixty{} to \maeXgbPerSegmentSixty{} mph.
Neighbour features give a small but statistically consistent gain over the pooled model (\gapAllFifteen{} mph at 15 min), concentrated in peak hours and congested states, but the gain does not grow with horizon.
This motivates, but does not prove, the benefit of learned graph models studied in Phase 2.
\end{abstract}

\section{Introduction}
Congestion on a freeway does not stay where it starts: a jam's tail moves upstream, so the state of downstream neighbours is informative about a sensor's near future.
Per-segment models see only their own series and can use this only indirectly.
Our hypothesis is that graph-aware models help most at longer horizons and during peaks.
Phase 1 builds the data pipeline and evaluation protocol, tests how far segment-level machine learning goes, and quantifies, with a ``bridge'' experiment, how much of the graph benefit hand-made spatial features capture.
Contributions: (i) a leakage-safe protocol with masked metrics; (ii) five ML forecasters and three baselines evaluated on all horizons, peak/off-peak and congested/free-flow slices; (iii) a gap analysis with paired block-bootstrap confidence intervals.

\section{Related Work}
Classical approaches use historical averages and ARIMA-type models.
Gradient-boosted trees \cite{chen2016xgboost} and random forests \cite{breiman2001rf} are strong tabular baselines.
Graph neural networks model the road network explicitly: DCRNN \cite{li2018dcrnn} uses diffusion convolution in a recurrent encoder--decoder (and introduced the METR-LA preprocessing used here), STGCN \cite{yu2018stgcn}, Graph WaveNet \cite{wu2019graphwavenet}, ASTGCN \cite{guo2019astgcn} and GMAN \cite{zheng2020gman} add convolutions, adaptive graphs and attention.
Graph attention \cite{velickovic2018gat,brody2022gatv2} is the basis of Phase 2.

\section{Dataset and EDA}
METR-LA \cite{li2018dcrnn,jagadish2014bigdata} contains \graphNodes{} sensors at 5-minute resolution from \splitTrainStart{} to \splitTestEnd{} (34,272 timestamps). It records \textbf{speed only}; there is no flow or occupancy.
A reading of 0 mph denotes a sensor outage, not stopped traffic; \missingRate\% of readings are 0, and \networkOutages{} timestamps have at least 90\% of sensors missing simultaneously.
We therefore mask zeros in losses and metrics and never impute targets.
Mean speed is \meanSpeed{} mph; \congPeak\% of weekday-peak observations are below 40 mph versus \congOffpeak\% off-peak (Fig.~\ref{fig:eda}).
The road graph has \graphEdges{} directed edges; its largest strongly connected component has \graphSCC{} sensors (diameter \graphDiameter{}).
Median pairwise correlation falls with graph distance, from \corrOneHop{} at one hop to \corrTwoHop{} at two and \corrEightHop{} at eight hops, so neighbours are informative.
Directional evidence is weaker: among \xcorrPairs{} adjacent pairs, the cross-correlation peak is at a positive lag (upstream sensor reacts after downstream) for \xcorrAllPos\% of pairs overall and \xcorrPeakPos\% at peak hours, versus \xcorrPeakNeg\% negative at peak hours.
This supports spatial dependence clearly and upstream-moving propagation only modestly.
The test period has a higher outage rate (\missingTest\%) than training (\missingTrain\%).

\begin{figure}[t]\centering
\includegraphics[width=\columnwidth]{02_daily_profile.pdf}
\caption{Mean speed by time of day (band: IQR across sensors); shaded regions are the peak windows.}\label{fig:eda}
\end{figure}

\section{Problem Formulation}
Let $x_t\in\mathbb{R}^{N}$ be the speeds at step $t$ and $m_t\in\{0,1\}^N$ the observation mask.
Given history $X_{t-H+1:t}$ ($H=12$, one hour) the task is to predict $y_{t+h}$ for $h=1..F$ ($F=12$).
Metrics are masked MAE, RMSE and MAPE in mph, e.g.\ $\mathrm{MAE}=\sum m|y-\hat y|/\sum m$.
MAPE uses only targets of at least 1~mph.

\section{Feature Engineering}
For each (sensor, anchor time) we use: 12 lagged speeds; rolling mean/std/min/max over 3, 6, 12 steps; first differences; missing flags and counts; cyclic time-of-day and day-of-week encodings (so 23:55 and 00:00 are close), weekend and holiday flags; for pooled models also sensor id (categorical), latitude, longitude and in/out-degree.
XGB-spatial adds adjacency-weighted neighbour speeds (upstream and downstream separately, lags 0--2), 2-hop means, and the difference between own and neighbour speed.
Inputs are imputed (forward fill up to 30 minutes, then train-only time-of-day means) and flagged.
The daily lag $t-288$ lies outside the history window and is kept off in the main comparison.

\section{Methods}
\textbf{Baselines.} Persistence, historical average (per day-of-week and time-of-day, train only), seasonal naive (same time yesterday).
\textbf{XGBoost.} Additive trees with a second-order objective and regularisation $\gamma T+\tfrac12\lambda\|w\|^2$ \cite{chen2016xgboost}; Optuna chooses the loss (absolute error was selected) and tree parameters on the validation split.
\emph{Per-segment}: one model per sensor and horizon (15, 30, 60~min); \emph{pooled}: one model per horizon across sensors.
\textbf{Random Forest.} Bagged deep trees with feature subsampling \cite{breiman2001rf}, natively multi-output (one forest predicts all 12 horizons).
\textbf{XGB-spatial.} The pooled model plus the neighbour features above.

\section{Experimental Setup}
Splits are chronological by timestamp: \splitTrainStart--\splitTrainEnd{} train, \splitValStart--\splitValEnd{} validation, \splitTestStart--\splitTestEnd{} test.
Windows are built inside each split so no input or target crosses a boundary.
Scalers, imputation means, the historical-average table and all hyperparameters use train (or validation for tuning) only.
Tuning budget: 15 Optuna trials per XGBoost variant and 6 for the forest (reduced to suit a laptop); pooled models are trained on a random 25\% of training rows and the forest on 10\%, which likely understates their accuracy.
Models are retrained on train only for test evaluation.
Significance uses the Diebold--Mariano test \cite{diebold1995comparing} with a Newey--West long-run variance \cite{newey1987simple} and the small-sample correction of \cite{harvey1997testing}, and day-block bootstrap intervals.

\section{Results}
Table~\ref{tab:main} gives test errors; Fig.~\ref{fig:horizon} plots error versus horizon.
All learned models beat persistence at 15 minutes (\maePersistenceFifteen{} mph); the historical average is flat across horizons (\maeHistoricalAverageFifteen{} mph) because it ignores the current state.
XGB-spatial is best at 15 and 30 minutes (\maeXgbSpatialFifteen{} and \maeXgbSpatialThirty{} mph) while per-segment XGBoost is best at 60 minutes (\maeXgbPerSegmentSixty{} mph vs.\ \maeXgbSpatialSixty{}); the forest is weakest of the ML models (\maeRfMultioutputFifteen{}/\maeRfMultioutputThirty{}/\maeRfMultioutputSixty{} mph).
No result violates the sanity bound for leakage (MAE at 60 minutes below about 3.0 mph).

\begin{table}[t]
\centering
\caption{Test-set error (mph) by model and horizon (minutes).}\label{tab:main}
\footnotesize\resizebox{\columnwidth}{!}{%
\begin{tabular}{lrrrrrr}
\toprule
model & MAE@15 & MAE@30 & MAE@60 & RMSE@15 & RMSE@30 & RMSE@60 \\
\midrule
historical average & 4.19 & 4.19 & 4.19 & 7.86 & 7.86 & 7.86 \\
persistence & 3.49 & 4.20 & 5.37 & 6.39 & 8.02 & 10.30 \\
rf multioutput & 3.14 & 3.80 & 4.73 & 5.70 & 6.89 & 8.25 \\
seasonal naive & 5.50 & 5.50 & 5.50 & 10.71 & 10.71 & 10.71 \\
xgb global & 2.88 & 3.33 & 3.78 & 5.69 & 6.80 & 7.76 \\
xgb per segment & 2.86 & 3.28 & \textbf{3.68} & 5.60 & 6.62 & \textbf{7.52} \\
xgb spatial & \textbf{2.82} & \textbf{3.26} & 3.73 & \textbf{5.46} & \textbf{6.54} & 7.54 \\
\bottomrule
\end{tabular}}
\end{table}


\begin{figure}[t]\centering
\includegraphics[width=\columnwidth]{error_vs_horizon.pdf}
\caption{Test MAE versus horizon for baselines and Phase 1 models.}\label{fig:horizon}
\end{figure}

\textbf{Feature importance.} TreeSHAP on 5,000 test rows shows that short-window averages and the latest speed dominate at 15 minutes, while sensor identity and time of day gain weight at 60 minutes; neighbour speeds appear in the top features of XGB-spatial but never dominate.

\textbf{Gap analysis.} Fig.~\ref{fig:gap} shows MAE(pooled)$-$MAE(spatial) by horizon and slice.
At 15 minutes the spatial features reduce MAE by \gapAllFifteen{} mph overall (95\% CI \gapAllFifteenLo--\gapAllFifteenHi), by \gapPeakFifteen{} mph at peak versus \gapOffpeakFifteen{} off-peak, and by \gapCongestedFifteen{} mph in congested states.
At 60 minutes the gains are \gapAllSixty{} (all), \gapPeakSixty{} (peak), \gapOffpeakSixty{} (off-peak) and \gapCongestedSixty{} (congested) mph.
The benefit therefore concentrates in peaks and congestion, as the propagation hypothesis predicts, but it does not increase with horizon.

\begin{figure}[t]\centering
\includegraphics[width=\columnwidth]{gap_analysis.pdf}
\caption{Gap analysis with paired block-bootstrap 95\% intervals (positive: spatial features help).}\label{fig:gap}
\end{figure}

\section{Discussion and Limitations}
Hand-made neighbour features capture part of the spatial signal cheaply; whether a learned graph model captures more is the Phase 2 question.
Limitations: one city and season; speed-only data; a fixed distance-based graph; reduced tuning and subsampled training rows; MAPE is unstable at low speeds; the test period has more outages than training.

\section{Conclusion and Phase 2 Plan}
Segment-level ML is a strong baseline and neighbour information helps modestly, mainly in congestion.
Phase 2 adds a spatio-temporal graph attention network, a no-graph deep ablation, a residual hybrid with leakage-safe out-of-fold base forecasts, ablations, maps and a simulated control experiment.

\bibliographystyle{IEEEtran}
\bibliography{../refs}
\end{document}
